% =============================================================================
% Appendices
% Only the slots relevant to this empirical study are kept. Consent forms and
% questionnaires are not included because the project has no human
% participants (see project_ethics.tex).
% =============================================================================
\appendix

\chapter{Installation and Reproduction Guide}
\label{app:install}

The full reproduction pipeline used to generate every figure
and table in this dissertation is documented at the project
repository root. The short version is:

\begin{code}[language=bash, caption={End-to-end reproduction
    commands.}, label={lst:reproduce}]
git clone https://github.com/[your-handle]/COOM.git
cd COOM
conda env create -f environment.yml   # or follow setup.py
conda activate coom

# Run all five experiments sequentially (~10 hours on Apple-Silicon CPU)
bash CL/scripts/cpu_queue.sh

# Regenerate figures and tables
PYTHONPATH=. python paper/scripts/make_figures.py
\end{code}

The reproduction has three external dependencies: a working
ViZDoom binary (shipped via PyPI, installed automatically by
\texttt{pip install .}), the OpenGL stack provided by macOS
itself, and roughly $10$~hours of wall-clock time on an
M-series CPU. No GPU is required and no cloud account is
required.

The figures and tables produced by
\texttt{paper/scripts/make\_figures.py} are written into the
\texttt{paper/figs/} and the project root, with names that match
the \texttt{\textbackslash input} commands in the corresponding
chapters. Re-running the script does not invalidate any
previously rendered PDF, because both the figure files and the
TSV inputs are content-addressed by the run identifier.

\chapter{Selected Code Listings}
\label{app:code}

The two listings below are reproduced because both are
load-bearing for the dissertation: the first for the analysis
pipeline that satisfies \ref{req:reproduce}, the second for the
overnight queue infrastructure described in \Cref{sec:queue}.

\begin{code}[language=Python, caption={Loader for COOM TSV
    training logs, handling the trailing-tab schema growth that
    arises when the upstream code appends new metric columns
    mid-run.}, label={lst:loader}]
def load_run(group_dir):
    path = latest_progress(group_dir)
    with path.open() as f:
        header = f.readline().rstrip('\n').split('\t')
    names = header + [f"_pad{i}" for i in range(20)]
    df = pd.read_csv(path, sep='\t', skiprows=1, names=names,
                     engine='c', index_col=False)
    return df[header]
\end{code}

\begin{code}[language=bash, caption={Full overnight CPU
    experiment queue, recording start time, return code and
    duration to a status file.}, label={lst:queuefull}]
#!/usr/bin/env bash
set -u
STATUS=CL/logs/queue_status.txt
LOG=CL/logs/queue_run.log

run() {
    local tag="$1"; shift
    local t0=$(date +%s)
    echo "[$(date '+%F %T')] START $tag :: $*" | tee -a "$STATUS"
    "$@" >> "$LOG" 2>&1
    local rc=$?
    local t1=$(date +%s)
    echo "[$(date '+%F %T')] DONE  $tag :: rc=$rc  duration=$((t1 - t0))s" \
        | tee -a "$STATUS"
}

run ft_seed1   python CL/run_cl.py --sequence CO4 --seed 1 \
    --steps_per_env 20000 --no_test --group_id ft_co4_20k_seed1
run l2_seed1   python CL/run_cl.py --sequence CO4 --seed 1 \
    --steps_per_env 20000 --no_test \
    --cl_method l2 --cl_reg_coef 100000 --group_id l2_co4_20k_seed1
run mas_seed1  python CL/run_cl.py --sequence CO4 --seed 1 \
    --steps_per_env 20000 --no_test \
    --cl_method mas --cl_reg_coef 10000 --group_id mas_co4_20k_seed1
run ewc_seed2  python CL/run_cl.py --sequence CO4 --seed 2 \
    --steps_per_env 20000 --no_test \
    --cl_method ewc --cl_reg_coef 250 --group_id ewc_co4_20k_seed2
run ewc_seed3  python CL/run_cl.py --sequence CO4 --seed 3 \
    --steps_per_env 20000 --no_test \
    --cl_method ewc --cl_reg_coef 250 --group_id ewc_co4_20k_seed3

echo "[$(date '+%F %T')] QUEUE COMPLETE" | tee -a "$STATUS"
\end{code}
